\documentclass{article}
\usepackage{graphicx} % Required for inserting images
\usepackage[utf8]{inputenc} % Per caratteri accentati
\usepackage[T1]{fontenc}    % Per una corretta codifica dei font
\usepackage[italian]{babel} % Per sillabazione in italiano
\usepackage{float}
\usepackage{hyperref}
\usepackage{listings}
\usepackage{color}
\usepackage{amssymb}
\usepackage{amsmath}
\usepackage{multicol}

%VALORI SISTEMABILI
\usepackage[a4paper, top= 1 cm, bottom= 1.5 cm, left= 2 cm, right= 2 cm]{geometry}

\title{Linguaggi di programmazione}
\author{Francesco Lanzarini}
\date{ }

\begin{document}

\maketitle

\section{Capitolo 1 - Macchine astratte, interpreti, compilatori}
\subsection{\textcolor{blue}{Che cosa è una macchina astratta?} E in cosa si differenzia da una macchina fisica?}
Una macchina astratta e' un insieme di strutture dati e algoritmi in grado di memorizzare e esegiure programmi, la parte esseziale di essa e' l'interprete.
Una MA si differenza da una macchina fisica perche' e' capace, tramite alcune fasi, di eseguire programmi scritti in linguagggi diversi; infatti una macchina
fisica e' in grado di eseguire programmi scritti nel suo linguaggio.
\subsection{\textcolor{red}{Che cos'è un interprete?} In cosa consiste il ciclo fetch-decode-execute?}
L'interprete e' il componente essenziale di una macchina astratta ed ha il compito di interpretare le istruzioni, esso e' costituito da: operazioni per
elaborare istruzioni primitive, operazioni e strutture dati per controllare la sequenza di esecuzione, trasferimento dati e gestione memoria.
Il ciclo fde consiste nel prelevare la giusta istruzione da eseguire (fetch), capire quale operzione e' necessario eseguire (decode) e eseguirla salvando eventualmente
il risultato in una certa zona di memoria (execute).
\subsection{Cos'è il linguaggio macchina?}
Il linguaggio macchina e' il linguaggio che la macchina e' in grado di comprendere ed eseguire.
\subsection{Possono esistere macchine diverse con lo stesso linguaggio macchina?}
Si, perche' la stessa istruzione elementare puo' essere implementata in modo differente in due macchine diverse.
\subsection{In quali modi è possibile implementare una macchina astratta? Elencare vantaggi e svantaggi delle varie tecniche.}
Sono tre i principali modi per realizzare una macchina astratta:
\begin{enumerate}
    \item Hardware: e' teoricamente sempre possibile ed ha la velocita' di esecuzione massima ma non e' per nulla flessbile per cui e' usata solo per macchine di basso livello o dedicate
    \item Microprogrammi in ROM: rende microprogrammabile la macchina ospite, ha una maggiore falessibilita' ma minore velocita'.
    \item Software di iterpretazione o simulazione: ovvero strutture dati e algoritmi della macchina astratta scritti in lunguaggio della macchina ospite, puo' essere implementata in qualsiasi
            macchina ospite, massima flessibilita' ma minima velocita'.
\end{enumerate}
Nella pratica in realta' si usa una combinazione di queste tecniche.
\subsection{\textcolor{red}{Che cos'è un compilatore?}}
Il compilatore e' un programma che traduce un programma da una lingua $L_1$ ad una lingua $L_0$ preservandone la semantica.
\subsection{Relativamente alla tecnica d'implementazione software, \textcolor{blue}{descrivere la tecnica d'implementazione interpretativa pura e quella compilativa pura}.}
Quando si esegue un'implementazione interpretativa pura si realizza un interprete scritto in $L_0$ (linguaggio macchina) per un linguaggio generico $L$, le singole istruzioni sono interpretate e poi eseguite.
Mentre in un'implementazione compilativa pura i programmi scritti in $L$ sono tradotti in programmi equivalenti in $L_0$ e poi eseguiti sulla macchina fisica

\subsection{\textcolor{red}{Quando un interprete si può dire corretto? Quando un compilatore si può dire corretto?}}
Un interprete si puo' dire corretto quando realizza una funzione parziale: $I^{L_0}_{L}: (P_r^{L} \times D)$  tale che $I_L^{L_0}(P_r^L, Input) = P^L(Input)$.
Un compilatore si puo' dire corretto quando realizza una funzione totale: $C_{L,L_0}:Prog^L \to Prog^L_0$ tale che $C_{L,L_0}(P_r^L)=P_rc^{L_0}$
\subsection{Confrontare l'implementazione di una macchina astratta su una macchina ospite per mezzo di un interprete o di un compilatore.}
In un'implementazione interpretativa pura si ha una scarsa efficienza della macchina astratta perche' bisogna sommare ai tempi di esecuzione i tempi di decodifica del codice,
inoltre nel caso una stessa istruzione dovesse ripetersi l'interprete deve decodificarla piu' volte; essa e' pero' un implementazione piu' flessibile, piu' facile da realizzare e non genera altro codice da memorizzare.
Mentre un'implementazione comilativa pura puo' vantare un'ottima velocita' visto che tutto il codice e' gia' tradotto, essa e' pero' piu' difficile da realizzare inoltre e' poco flesssibile
(bisognerebbe riscrivere il compilatore), astrae di piu' dal programma sorgente e occupa piu' memoria.
\subsection{Come vengono implementate nella realtà le macchine astratte? \textcolor{blue}{Che cos'è la macchina intermedia?}}
Nella realta' e un misto delle due interpretazioni: di solito si compila il programma in una lingua intermedia $L_i$ per una la macchina intermedia, suiccessivamente un interprete $I^{L_0}_{L_i}$ interpreta le istruzioni per l'esecuzione.
Generalmente $I^{L_0}_L_i$ = interprete della macchina ospite + opportuni meccanismi chiamati supporto a run tim (SRT).
\subsection{\textcolor{blue}{Quando si dice che una implementazione è di tipo interpretativo e quando di tipo compilativo?} Fare esempi di linguaggi la cui implementazione è di un tipo o dell'altro.}
Solitamente si dice che un implementazione e' interpretativa quando la maggior parte delle istruzioni vengono interpretate e poi eseguite, tipo in python e Java; mentre 
un'implementazione e' di tipo compilativo quando la maggior parte delle istruzioni sono prima tradotte e poi eseguite, tipo in C e C++.
\subsection{Cos'è la Just-in-Time compilation?}
La JIT e' una combinazione di AOT (ahead of time) compilation e interpretazione: si genera prima un codice intermedio (byte code in java) che viene poi eseguito passando per un compilatore
on-demand man mano che esegue il programma tradotto in linguaggio macchina.
\subsection{L'interprete e il compilatore si possono sempre realizzare?}
\subsection{Che cosa è l'implementazione via kernel?}
\subsection{Cosa si intende per bootstrapping?}

\section{Capitolo 2 - Descrivere un linguaggio di programmazione}
\subsection{Quali sono i livelli di descrizione di un linguaggio?}
\subsection{Sintassi: qual è l'aspetto lessicale e quale quello grammaticale di un linguaggio?}
\subsection{Cos'è un alfabeto? Cos'è una parola o stringa? Cos'è A*? Tale insieme è enumerabile?}
\subsection{Definizione di potenza di una stringa. Definizione di potenza di un linguaggio. Definizione di chiusura/iterazione (o stella di Kleene) di un linguaggio.}
\subsection{Definizione di grammatica libera da contesto. Come si deriva una stringa? Qual è il linguaggio generato da una grammatica libera?}
\subsection{Cos'è un albero di derivazione? Cos'è una derivazione canonica sinistra/destra (leftmost/rightmost)? Esiste una corrispondenza biunivoca tra alberi di derivazioni e derivazioni canoniche?}
\subsection{\textbf{Quando una grammatica è ambigua? (fare un esempio) Quando un linguaggio è ambiguo? (fare un esempio)}}
\subsection{È possibile rimuovere l'ambiguità dalla grammatica delle espressioni aritmetiche? Come?}
\subsection{Cos'è l'albero di sintassi astratta? Che differenza c'è tra sintassi concreta e sintassi astratta? Cos'è lo zucchero sintattico?}
\subsection{Fare esempi di vincoli sintattici contestuali. Possono essere catturati attraverso grammatiche libere?}
\subsection{\textit{Cosa s'intende per semantica statica? E per semantica dinamica?}}
\subsection{\textbf{Elencare le varie fasi in cui si articola un compilatore. Descrivere in dettaglio ogni singola fase.}}
\subsection{A chi serve definire la semantica dinamica di un linguaggio e perché?}
\subsection{Quali tecniche si usano per dare semantica dinamica ad un linguaggio di programmazione?}
\subsection{Imparare le regole di semantica operazionale SOS per il semplice linguaggio presentato a lezione. Regole di valutazione: interna vs esterna, sinistra vs destra vs parallela.}
\subsection{Cosa s'intende per pragmatica? Cosa si intende per implementazione di un linguaggio?}

\section{Capitolo 3 - Analisi lessicale-Linguaggi regolari}
\subsection{Cosa fa l'analizzatore lessicale?}
\subsection{Cos'è un token?}
\subsection{Cos'è un pattern? Come lo si rappresenta? Cos'è un lessema?}
\subsection{\textbf{Definizione di espressioni regolari (sintassi) e di linguaggio associato (semantica).}}
\subsection{Quali sono i linguaggi regolari? I linguaggi finiti sono tutti regolari? Esistono linguaggi infiniti regolari?}
\subsection{Definizione di equivalenza tra espressioni regolari. Elencare alcune leggi di equivalenza.}
\subsection{Cos'è una definizione regolare e a cosa serve?}
\subsection{\textbf{Definizione di NFA (automa finito nondeterministico). Discutere cosa si intende per nondeterminismo. Mettere in relazione la definizione formale con la rappresentazione grafica come diagramma di transizioni.}}
\subsection{Come si definisce il linguaggio accettato da un NFA? Definizione di equivalenza tra NFA.}
\subsection{\textbf{Definizione di DFA (automa finito deterministico). Discutere cosa si intende per determinismo. Mostrare che un DFA è un caso speciale di NFA, ovvero la classe dei DFA è un sottoinsieme della classe degli NFA.}}
\subsection{\textbf{Dato un NFA, come si ricava un equivalente DFA? Ovvero descrivere come è definita la costruzione per sottoinsiemi. Definizione di epsilon-closure e algoritmo associato. Qual è la complessità della costruzione per sottoinsiemi nel caso pessimo? Ovvero se NFA ha n stati, quanti stati può avere il DFA equivalente?}}
\subsection{\textbf{Dato i due punti precedenti, enunciare il teorema che dice che la classe dei linguaggi riconosciuto da NFA coincide con la classe dei linguaggi riconosciuti da DFA.}}
\subsection{\textbf{Come si costruisce un NFA a partire da una espressione regolare, in modo tale che il linguaggio riconosciuto dall'NFA sia lo stesso del linguaggio associato all'espressione regolare?}}
\subsection{Definizione di grammatica regolare.}
\subsection{Come si associa ad una grammatica regolare un equivalente NFA?}
\subsection{Dato un DFA, come si costruisce una grammatica regolare equivalente?}
\subsection{Data una grammatica regolare, come si costruisce un'espressione regolare equivalente?}
\subsection{\textit{Descrivere, con un diagramma riassuntivo, tutte le relazioni fra i formalismi introdotti: NFA, DFA, grammatiche regolari, espressioni regolari. Questo diagramma dimostra che tutti questi formalismi sono equivalenti e descrivono la classe dei linguaggi regolari.}}
\subsection{Definizione di stati equivalenti in un DFA. Quando due stati di un DFA sono distinguibili?}
\subsection{\textit{Come funziona l'algoritmo iterativo con tabella a scala per produrre le classi di equivalenza di stati di un DFA?}}
\subsection{Una volta determinate le classi di equivalenza degli stati di un DFA, come si costruisce l'automa minimo associato?}
\subsection{\textit{Cos'è Lex? Qual è il suo input e il suo output?}}
\subsection{\textit{Qual è la struttura di un file .l? Come funziona l'analizzatore lessicale prodotto da Lex?}}
\subsection{Come si interfaccia Lex con Yacc?}
\subsection{\textbf{Intestazione e dimostrazione del pumping lemma.}}
\subsection{\textbf{Come si può utilizzare il pumping lemma (a rovescio) per dimostrare che un linguaggio non è regolare?}}
\subsection{\textit{Quali sono le proprietà di chiusura dei linguaggi regolari? Quali proprietà si possono decidere (ovvero verificare algoritmicamente)?}}

\section{Capitolo 4 - Analisi sintattica - linguaggi liberi}
\subsection{Cosa si intende per analisi sintattica? Cos'è un parser?}
\subsection{\textbf{Definizione di automa a pila nondeterministico (PDA). Definizione di configurazione (o descrizione istantanea), mossa in un passo e mossa in più passi.}}
\subsection{\textbf{Definizione di linguaggio accettato da un PDA per stato finale o per pila vuota. Sono equivalenti queste due modalità di riconoscimento?}}
\subsection{\textbf{Mostrare come data una grammatica libera G, sia possibile costruire un PDA P equivalente. È possibile costruire, dato un PDA P, una grammatica equivalente G? Concludere che la classe dei linguaggi liberi coincide con la classe dei linguaggi riconosciuti da PDA.}}
\subsection{\textit{Quali sono le proprietà di chiusura dei linguaggi liberi? L'intersezione di un linguaggio libero con un regolare è un linguaggio libero?}}
\subsection{\textbf{Intestazione e dimostrazione del "Pumping Theorem".}}
\subsection{\textbf{Come si può utilizzare tale teorema (a rovescio) per dimostrare che un linguaggio non è libero?}}
\subsection{\textit{Classificazione di Chomsky delle grammatiche e dei linguaggi. Definizione di grammatica dipendente dal contesto e di grammatica monotona. Quale tipo di automi corrisponde ad ogni classe?}}
\subsection{Definizione di DPDA (automa a pila deterministico) e di linguaggio libero deterministico.}
\subsection{La classe dei linguaggi liberi deterministici è strettamente inclusa in quella dei linguaggi liberi? Contiene strettamente la classe dei linguaggi regolari?}
\subsection{Che cosa dice la prefix property e perché è interessante per i DPDA?}
\subsection{Usando un endmarker \$, si può riconoscere un linguaggio libero deterministico che non gode della prefix property anche per pila vuota? Come?}
\subsection{Un linguaggio libero deterministico è ambiguo?}
\subsection{Proprietà di chiusura dei linguaggi liberi deterministici: chiusi per complementazione, ma non per intersezione né per unione.}
\subsection{Da cosa si parte per costruire un analizzatore sintattico (ovvero parser)? Da una espressione regolare? Da una grammatica libera? Da un PDA?}
\subsection{Cosa prende in input e cosa produce in output un parser?}
\subsection{Che differenza c'è tra un parser nondeterministico ed uno deterministico?}
\subsection{Quali sono le due tecniche essenziali per costruire parsers?}
\subsection{Le tecniche top-down e bottom-up in che cosa differiscono?}
\subsection{Quali tipi di grammatiche non sono adatte al top-down parsing? Quali tipi di produzioni sono poco adatte al bottom-up parsing?}
\subsection{Cosa sono le produzioni epsilon e cosa sono i simboli non-terminali annullabili?}
\subsection{\textit{Come si può trasformare una grammatica G che contiene produzioni epsilon in una grammatica G' che non ne contiene, preservando il linguaggio a meno di epsilon?}}
\subsection{Cosa sono le produzioni unitarie e cosa sono le coppie unitarie?}
\subsection{\textit{Come si può trasformare una grammatica G che contiene produzioni unitarie in una equivalente G' che non ne contiene?}}
\subsection{Data una grammatica G, quali sono i suoi simboli utili? Quali sono i suoi simboli generatori? Quali i suoi simboli raggiungibili?}
\subsection{\textit{Come si calcolano i generatori? Come si calcolano i raggiungibili?}}
\subsection{\textit{In che modo si eliminano i simboli inutili di una grammatica? È importante l'ordine delle operazioni da svolgere?}}
\subsection{\textbf{Quando si dice che una grammatica è ricorsiva sinistra? Come si elimina la ricorsione sinistra immediata? E quella non immediata? Perché serve eliminare la ricorsione sinistra?}}
\subsection{\textbf{Cosa vuol dire fattorizzare a sinistra una grammatica? Perché serve fattorizzare?}}
\subsection{\textbf{Definizione di First(a) e di Follow(A).}}
\subsection{\textbf{Algoritmi per calcolare First(a) e di Follow(A).}}
\subsection{\textbf{Come è fatta e come si riempie una tabella di parsing LL(1)?}}
\subsection{\textbf{Quando una grammatica G si dice di classe LL(1)? Quali sono le condizioni necessarie e sufficienti per G affinché sia di classe LL(1)?}}
\subsection{Perché un parser di questo tipo è chiamato LL?}
\subsection{\textbf{Come funziona il parser LL(1) con una pila?}}
\subsection{\textit{Cos'è un parser a discesa ricorsiva?}}
\subsection{È vero che ogni linguaggio regolare è pure LL(1)?}
\subsection{Come sono definiti Firstk(a) e di Followk (A) per $k\ge2$.}
\subsection{Come si definisce una grammatica LL(k)? Ed un linguaggio LL(k)? Come si relazionano tra di loro?}
\subsection{Che relazione esiste tra grammatiche LL(k) e grammatiche ambigue? E con le grammatiche ricorsive sinistre?}
\subsection{Esistono linguaggi liberi che non sono LL(k) per nessun k? Ed esistono linguaggi liberi deterministici che non sono LL(k) per nessun k?}
\subsection{\textbf{Cos'è un parser bottom-up (o shift-reduce)? Qual è il suo input e il suo output? Perché sono chiamati parser LR?}}
\subsection{\textbf{Che tipo di conflitti si possono presentare in un parser del genere? Quando si presenta un conflitto (shift/reduce o reduce/reduce), quale azione bisogna scegliere?}}
\subsection{\textit{Cos'è un prefisso viabile? Come lo si definisce in termini di una grammatica?}}
\subsection{\textit{Cos'è un item LR(0)? Come si generano tutti gli item di una grammatica (aumentata con un simbolo iniziale nuovo S')?}}
\subsection{\textbf{Come si ricava il DFA dei prefissi viabili, detto anche automa canonico LR(0)?}}
\subsection{\textbf{Come è fatta una tabella di parsing LR(0)? Come la si riempie a partire dall'automa canonico LR(0)? Quando una grammatica è di classe LR(0)?}}
\subsection{\textbf{Come è fatto il parser LR(0) che utilizza la tabella di parsing LR(0)? Quanti stack servono?}}
\subsection{Esistono grammatiche non LR(0)? Fare un esempio semplice.}
\subsection{\textbf{Come è fatta e come si riempie una tabella di parsing SLR(1)? Cosa vuol dire l'acronimo SLR? Perché si mette 1 come parametro?}}
\subsection{\textbf{Quando una grammatica è di classe SLR(1)? Esistono grammatiche non di classe SLR(1)?}}
\subsection{\textit{Cos'è un item LR(1)? Come si costruisce l'automa canonico LR(1)?}}
\subsection{\textit{Come è fatta e come si riempie una tabella di parsing LR(1)?}}
\subsection{\textit{Come è fatto una tabella di parsing LALR(1)? Quando una grammatica è di classe LALR(1)?}}
\subsection{Esistono grammatiche LALR(1) che non sono SLR(1)? Esistono grammatiche LR(1) che non sono LALR(1)? Mostrare una grammatica che è LR(1) ma non LALR(1).}
\subsection{Come si può generalizzare l'idea per ogni $k \geq 2$? Ovvero quando una grammatica è di classe LR(k), SLR(k) o LALR(k)? }
\subsection{Come si relazionano le grammatiche della famiglia LR (LR, SLR, LALR) al variare di k? Come si relazionano le grammatiche LR(k) e LL(k) al variare di k?}
\subsection{Le grammatiche LR(k) e LL(k) sono sempre non ambigue? Esistono grammatiche ambigue che sono LL(k) o LR(k) per qualche k? Esistono grammatiche che non sono LR(k) per nessun k?}
\subsection{Come si relazionano i linguaggi della famiglia LR(k) rispetto a quelli LL(k)? Esistono linguaggi liberi deterministici che non sono LR(k) per qualche k? Esistono linguaggi liberi deterministici che non sono LL(k) per qualche k?}
\subsection{Esiste un linguaggio regolare che non è LR(0)? Come si relazionano i linguaggi LR(0) rispetto a quelli LL(1)? La prefix property è una condizione necessaria e sufficiente affinché un linguaggio sia di classe LR(0)?}
\subsection{La classe dei linguaggi SLR(1) coincide con la classe dei linguaggi liberi deterministici?}
\subsection{\textit{Cos'è YACC? Qual è il suo input e il suo output? Come si ottiene un parser eseguibile a partire da un file .y? Come agisce il parser generato da YACC in sintonia con lo scanner generato da Lex?}}
\subsection{\textit{Come è la struttura di un file .y di YACC? Cosa sono le regole? Cos'è l'azione semantica?}}
\subsection{\textit{È possibile gestire grammatiche ambigue con YACC, specificando le associatività e le priorità fra gli operatori per risolvere l'ambiguità? Come si comporta YACC in presenza di conflitti?}}

\section{Capitolo 5 - Fondamenti}
\subsection{\textbf{È possibile costruire un programma Check che, preso in input un qualunque programma P, restituisce 1 se P è corretto e 0 se P è scorretto? Ovvero existe un qualche compilatore che può scovare tutti i possibili errori di un programma?}}
\subsection{\textbf{Cosa dice il problema della fermata (Halting Problem)? (L'errore in esame è la possibilità di non terminare il calcolo.) Come si dimostra che il problema non può essere risolto?}}
\subsection{\textit{Quando un problema è decidibile?}}
\subsection{\textit{Quali sono tipici esempi di proprietà indecidibili per i linguaggi di programmazione?}}
\subsection{\textbf{Cos'è una Macchina di Turing (MdT)? Che cosa calcola?}}
\subsection{\textbf{Quando un linguaggio è detto Turing-completo? Cosa afferma la tesi di Church-Turing e perché non si può dimostrare?}}
\subsection{I normali linguaggi di programmazione sono Turing-completi? Cosa afferma il teorema di Böhm-Jacopini?}
\subsection{\textit{Quale relazione existe tra espressività di un formalismo e decidibilità di proprietà dello stesso? Fare una panoramica prendendo in esame i tre formalismi MdT, PDA, DFA, e le due proprietà w$\in$L(M) (è w riconosciuta dalla macchina M?) e $L(M_1)=L(M_2)$ (le due macchine sono equivalenti, ovvero riconoscono lo stesso linguaggio?).}}



\end{document}